% Literature review: Fractional Reservoir Computing for Critical Brain Dynamics
% Companion to the fracres library (https://github.com/dave2k77/fracres).
% Build: tectonic main.tex   (or upload the folder to Overleaf)
\documentclass[11pt,a4paper]{article}

\usepackage[margin=2.5cm]{geometry}
\usepackage{amsmath,amssymb,amsthm}
\usepackage{booktabs}
\usepackage[colorlinks=true,linkcolor=blue,citecolor=blue,urlcolor=blue]{hyperref}
\usepackage[numbers,sort&compress]{natbib}
\usepackage{microtype}

\newtheorem{theorem}{Theorem}[section]
\newtheorem{definition}[theorem]{Definition}
\newtheorem{remark}[theorem]{Remark}

\title{Fractional Reservoir Computing for Critical Brain Dynamics:\\
       A Review of the Mathematical and Modelling Foundations}
\author{D.~R.~Chin\\\small Companion document to the \texttt{fracres} library}
\date{September 2026}

\begin{document}
\maketitle

\begin{abstract}
\noindent
The \texttt{fracres} library builds reservoir computers whose nodes carry
fractional (power-law, non-Markovian) memory, driven by fractional Gaussian
noise and regularised toward critical dynamics in a Besov/$L^p$ geometry, with
the aim of generating and analysing synthetic macroscopic neural signals with
realistic long-range dependence (LRD) and critical avalanche statistics. This
review surveys the four bodies of theory the construction rests on ---
(i)~fractional calculus and the stability of fractional-order systems,
(ii)~reservoir computing and the echo-state/edge-of-chaos literature,
(iii)~neural-mass and neural-field models and their fractional
generalisations, and
(iv)~critical brain dynamics, neuronal avalanches, self-organised and
quasi-self-organised criticality with homeostatic control --- together with
the mathematics of long-range dependence: self-similar processes, heavy tails,
and wavelet/Besov estimation. For each area we state the definitions and
theorems the library uses, summarise the neurophysiological evidence, and
identify the open gaps the \texttt{fracres} programme is positioned to probe.
All references were verified programmatically against publisher/Crossref,
OpenAlex, or arXiv records in September 2026.
\end{abstract}

\tableofcontents
\newpage

\section{Introduction and scope}
\label{sec:intro}

Macroscopic neural signals --- EEG, MEG, local field potentials --- exhibit two
statistical signatures that any generative model of brain dynamics must
confront: \emph{long-range dependence} (power-law temporal correlations and
$1/f$-type spectra persisting over hundreds of seconds
\cite{linkenkaer2001long,he2014scale}) and \emph{criticality} (power-law
distributions of neuronal avalanche sizes and durations, with exponents near
the mean-field branching values $\tau \approx 1.5$ and $\alpha \approx 2.0$
\cite{beggs2003neuronal,chialvo2010emergent}). Classical generative models
struggle with both: reservoir computers forget exponentially
\cite{jaeger2002memory}, neural-mass models relax exponentially
\cite{jansen1995electroencephalogram}, and strict self-organised criticality
requires conservation laws that brains do not obey \cite{bonachela2009self}.

The \texttt{fracres} project takes the position that these failures share a
common mathematical root: the dynamics are built from \emph{integer-order}
(exponential-memory) operators, while the physics of neural tissue --- from
single-neuron adaptation \cite{lundstrom2008fractional} to population
filtering --- is better described by \emph{fractional-order} (power-law
memory) operators. The library therefore places long memory in the
\emph{nodes} of a reservoir computer via discrete fractional derivatives,
drives the system with fractional Gaussian noise of tunable Hurst exponent,
stabilises it at the fractional edge of chaos via Matignon's stability
criterion, regulates it with a quasi-self-organised-criticality (qSOC)
homeostat, and constrains its outputs to the Besov-regularity manifold that
heavy-tailed long-memory signals inhabit. We call the resulting generative
models \emph{phantom brains}: synthetic neural signal generators with known,
controllable ground truth, intended both as hypothesis testbeds (which
mechanisms produce which statistics?) and as calibration phantoms for
estimators (which LRD/criticality estimator survives heavy tails and
structural breaks?).

This review assembles the theoretical foundations of that programme. Each
section states the mathematics the code implements, reviews the supporting
literature, and closes with the implications for the library. Citations were
collected and verified programmatically; residual verification caveats are
listed in Appendix~\ref{app:verification}.

\paragraph{Organisation.}
Section~\ref{sec:fractional}: fractional calculus, Mittag-Leffler relaxation,
Matignon stability, and numerical schemes.
Section~\ref{sec:rc}: reservoir computing, the echo state property, memory
capacity, and the small but growing fractional-RC literature.
Section~\ref{sec:masses}: neural-mass and neural-field models and their
fractional-order generalisations.
Section~\ref{sec:criticality}: neuronal avalanches, SOC, quasi-criticality,
homeostasis, and E/I balance.
Section~\ref{sec:lrd}: self-similar processes, LRD estimators, heavy tails,
and the neural evidence.
Section~\ref{sec:synthesis}: how the pieces assemble in \texttt{fracres}, and
the open research questions.

\section{Fractional calculus and fractional-order systems}
\label{sec:fractional}

\subsection{Definitions: Riemann--Liouville, Caputo, Gr\"unwald--Letnikov}

Fractional calculus extends differentiation and integration to non-integer
order; the modern standard references are Podlubny
\cite{podlubny1999fractional}, Kilbas, Srivastava \& Trujillo
\cite{kilbas2006theory}, and the classic treatment of Oldham \& Spanier
\cite{oldham1974fractional}. For $0 < \alpha < 1$ the two derivatives used in
\texttt{fracres} are the \emph{Caputo} derivative
\begin{equation}
  {}^C D^\alpha f(t) = \frac{1}{\Gamma(1-\alpha)}
  \int_0^t \frac{f'(\tau)}{(t-\tau)^{\alpha}}\,d\tau,
  \label{eq:caputo}
\end{equation}
which respects classical initial conditions $f(0)$, and the
\emph{Gr\"unwald--Letnikov} (GL) derivative, defined directly as a limit of
weighted differences,
\begin{equation}
  D^\alpha f(t) = \lim_{h \to 0} \frac{1}{h^\alpha}
  \sum_{j=0}^{\lfloor t/h \rfloor} c_j^{(\alpha)} f(t - jh),
  \qquad
  c_0 = 1,\quad c_j = \Bigl(1 - \frac{1+\alpha}{j}\Bigr) c_{j-1},
  \label{eq:gl}
\end{equation}
whose generalised-binomial weights $c_j^{(\alpha)}$ decay as a power law
$j^{-\alpha-1}$ \cite{oldham1974fractional,scherer2011grunwald}. The contrast
with integer order is the point: a first derivative is \emph{local} (two
samples), while a fractional derivative is a \emph{convolution over the whole
history} with a power-law kernel --- non-Markovian memory is definitional,
not engineered. For power laws both derivatives agree with the closed form
\begin{equation}
  D^\alpha t^\beta =
  \frac{\Gamma(\beta+1)}{\Gamma(\beta+1-\alpha)}\, t^{\beta-\alpha},
  \qquad \beta > -1,
  \label{eq:powerlaw}
\end{equation}
\cite[Eq.~2.117]{podlubny1999fractional}, which \texttt{fracres} uses as
ground truth to validate its discrete kernels.

\subsection{Mittag-Leffler relaxation}

The eigenfunction of the Caputo derivative is the Mittag-Leffler function
$E_\alpha(z) = \sum_{k\ge0} z^k/\Gamma(\alpha k + 1)$: $x(t) = E_\alpha(\lambda
t^\alpha)$ solves ${}^C D^\alpha x = \lambda x$, $x(0)=1$
\cite{podlubny1999fractional,gorenflo2014mittag}. Its asymptotics interpolate
between a fast (stretched-exponential) initial decay and a power-law tail
$x(t) \sim t^{-\alpha}$ \cite{gorenflo2014mittag,haubold2011mittag}: the
impulse response of a fractional-order system \emph{is} the non-exponential
fading-memory profile that neural LRD demands, in contrast to the exponential
forgetting of integer-order reservoirs (Section~\ref{sec:rc}). The
Caputo/Riemann--Liouville distinction matters in practice: with $x(0) \neq
0$, ${}^C D^\alpha f = D^\alpha_{\mathrm{RL}}(f - f(0))$, a correction the
library's validation suite exercises explicitly on the Mittag-Leffler
eigenfunction.

\subsection{Stability: Matignon's theorem}

For the commensurate linear fractional system $D^\alpha x = A x$,
$0<\alpha<1$, the stability criterion replaces the classical left-half-plane
condition by an \emph{argument} condition:

\begin{theorem}[Matignon \cite{matignon1996stability}]
  The system $D^\alpha x = A x$ is asymptotically stable if and only if
  every eigenvalue $\lambda_i$ of $A$ satisfies
  \begin{equation}
    |\arg(\lambda_i)| > \frac{\alpha \pi}{2}.
    \label{eq:matignon}
  \end{equation}
\end{theorem}

The unstable set is the cone of half-angle $\alpha\pi/2$ about the positive
real axis. As $\alpha \to 1$ it fills the right half-plane, recovering the
classical condition; for $\alpha < 1$ it \emph{shrinks}, so a
fractional-order system tolerates eigenvalues with positive real part --- a
strictly larger stable region \cite{matignon1996stability,matignon1998stability}.
This is the fractional analogue of the echo-state spectral condition
(Section~\ref{sec:rc}): it defines the wedge that the reservoir Jacobian
$A = W_{\mathrm{res}} - \lambda I$ must respect, and it is the rigorous sense
in which a fractional reservoir can sit closer to criticality --- with a
``hotter'' connectome --- than its integer-order counterpart. Matignon's 1998
extension covers non-commensurate (multi-order) systems
\cite{matignon1998stability}, the relevant theory once nodes are allowed
heterogeneous orders $\alpha_i$.

\subsection{Numerical schemes}

Two discretisations are implemented and validated in \texttt{fracres}. The
\emph{GL scheme}~\eqref{eq:gl} is first-order, $O(h)$
\cite{oldham1974fractional,scherer2011grunwald}; the \emph{L1 scheme} for the
Caputo derivative interpolates $f'$ piecewise-linearly on each subinterval,
giving weights $b_m = (m+1)^{1-\alpha} - m^{1-\alpha}$ and accuracy
$O(h^{2-\alpha})$ \cite{linxu2007l1}. The library's measured convergence
orders match both rates exactly (GL $\approx 1.00$; L1 $\approx 2-\alpha$ for
$\alpha \in \{0.3, 0.5, 0.7, 0.9\}$), confirming the implementations against
\eqref{eq:powerlaw}. For nonlinear node dynamics, the fractional
Adams--Bashforth--Moulton predictor--corrector of Diethelm, Ford \& Freed is
the standard higher-fidelity alternative \cite{diethelm2002predictor}; the
engineering literature also supplies rational (Oustaloup/CRONE) approximations
used to realise fractional operators in analogue and neuromorphic hardware
\cite{monje2010fractional,podlubny1999controllers}. All schemes face the same
practical issue: the full-history convolution is $O(TL)$, and truncation (the
short-memory principle \cite{scherer2011grunwald}) trades accuracy for cost
--- for hardware, the retained history at tolerance $\varepsilon$ scales as
$\varepsilon^{-1/\alpha}$ \cite{teuscher2026hardware}.

\subsection{Fractional-order reservoir computing}

The fractional-RC literature is small but direct. Yao \& Wang introduced the
first fractional-order echo state network, updating reservoir states through a
GL-type fractional difference and reporting improved prediction on chaotic
benchmarks \cite{yao2020fractional}. Goswami et al.\ prove the key theoretical
separation: standard ESNs induce \emph{short-memory} processes, while a
fractional-difference reservoir generates polynomially decaying autocorrelation
consistent with statistical long memory \cite{goswami2026longmemory} ---
precisely the property required to model LRD in neural signals. A 2026 review
positions fractional-order RC as a general extension of the RC paradigm,
covering the memory-capacity/fading-memory trade-off controlled by the order
$\alpha$ \cite{mastin2026fractional}, and Teuscher's hardware analysis
quantifies what physical fractional devices can deliver
\cite{teuscher2026hardware}. Notably absent from the literature: any
fractional reservoir targeted at \emph{critical} (avalanche) statistics, any
coupling of fractional nodes to neural-mass/field equations, and any
stochastic-drive (fGn) formulation --- the space \texttt{fracres} occupies.

\paragraph{Implications for \texttt{fracres}.}
The GL/L1 kernels (\texttt{kernels.py}) are validated against
\eqref{eq:powerlaw} and the Mittag-Leffler eigenfunction; Matignon's
criterion~\eqref{eq:matignon} is implemented as the edge-of-chaos control
(\texttt{stability.py}); the demonstrated ``fractional advantage'' --- one
connectome stable for $\alpha < 0.84$ where the $\alpha=1$ ESN diverges ---
is a direct numerical exhibition of Theorem~\ref{eq:matignon}.

\section{Reservoir computing}
\label{sec:rc}

\subsection{The reservoir/readout split}

Reservoir computing trains only a linear readout of a fixed, randomly wired
dynamical system. The two founding formulations are Jaeger's echo state
networks (ESNs) \cite{jaeger2001echo,jaeger2002adaptive} and Maass,
Natschl\"ager \& Markram's liquid state machines (LSMs)
\cite{maass2002liquid}, shown to be instances of one framework by Verstraeten
et al.\ \cite{verstraeten2007unification} and surveyed by Luko\v{s}evi\v{c}ius
\& Jaeger \cite{lukosevicius2009survey}. Two ideas from the foundations carry
directly into the present programme. First, the LSM's thesis that cortical
microcircuits compute on \emph{transients} rather than attractors
\cite{maass2002liquid} aligns with critical, avalanche-like dynamics.
Second, the essential object is the reservoir/readout split itself, not the
neuron model \cite{verstraeten2007unification} --- so replacing the node's
dynamical law (integer $\to$ fractional) while keeping the split is a
legitimate, minimal modification, which is what \texttt{fracres} does.

\subsection{Echo state property and fading memory}

A reservoir has the \emph{echo state property} (ESP) if its state is uniquely
determined by the left-infinite input history \cite{jaeger2001echo}; the
classical sufficient conditions bound the spectral radius (Jaeger) or, more
tightly, the largest singular value of $W_{\mathrm{res}}$
\cite{buehner2006tighter}. Yildiz, Jaeger \& Kiebel showed the ESP is
input-dependent: for some inputs no scaling of a given reservoir has echo
states \cite{yildiz2012echo}. Both theorems assume \emph{geometric}
contraction --- their Lipschitz bounds do not directly apply to power-law
memory, and re-establishing an ESP analogue under algebraic decay is an open
theoretical task for fractional RC. In \texttt{fracres} the role of the
contraction is played by the gain-normalised fractional update (unit linear
memory gain) plus the leak $-\lambda x$, with Matignon's wedge as the
linearised stability certificate; a formal fractional ESP proof remains
future work (\texttt{goswami2026longmemory} proves the complementary result
that fractional reservoirs \emph{do} produce long-memory stationary states
\cite{goswami2026longmemory}).

\subsection{Memory capacity}

Jaeger's memory-capacity curve measures how well the readout reconstructs
delayed inputs $u_{t-\tau}$; linear ESNs with $N$ units have total capacity
$\le N$, with recall decaying \emph{exponentially} in the delay
\cite{jaeger2002memory}. Dambre et al.\ generalised this to a full
information-processing capacity over all polynomial functionals of the input,
proving a conservation law: total capacity is bounded by the number of
independent reservoir variables \cite{dambre2012capacity}. The design
question for fractional RC is therefore not ``more capacity'' but
\emph{redistribution}: power-law kernels trade short-delay recall for
algebraically decaying long-delay recall within the same budget. The delayed-copy
benchmarks in \texttt{fracres} (test correlation $0.82 \to 0.51$ over delays
$1 \to 10$, and $0.64 \to 0.95$ after the $(h,\lambda)$ tuning study) measure
exactly this curve.

\subsection{Edge of chaos}

Bertschinger \& Natschl\"ager showed that real-time reservoir computation
peaks at the transition between ordered and chaotic dynamics --- the
\emph{edge of chaos} \cite{bertschinger2004edge}; Legenstein \& Maass extended
the analysis to biologically realistic circuits and proposed predicting
performance from the Jacobian spectrum \cite{legenstein2007edge}. In
classical RC the only route to the edge is tuning the spectral radius
$\rho(W_{\mathrm{res}}) \to 1$. A fractional reservoir has a second axis: the
order $\alpha$ moves the Matignon wedge~\eqref{eq:matignon}, so the
\emph{same} connectome can be unstable at $\alpha=1$ and stable at
$\alpha < 1$. Criticality tuning and memory tuning partially decouple --- a
structural advantage that the classical framework lacks.

\subsection{Modern context}

Physical and delay-based reservoirs replace topological memory with dynamical
memory in the substrate itself \cite{tanaka2019review,nakajima2021book} ---
the closest classical analogue of the fractional design. Next-generation
reservoir computing strips the reservoir to its memory-functional content
(nonlinear vector autoregression) \cite{gauthier2021nextgen}: both a
competitor and a natural host for fractional (long-memory) feature maps.

\paragraph{Implications for \texttt{fracres}.}
The library keeps the canonical split (frozen reservoir, trained linear
readout, ridge or gradient) and the delayed-copy memory probe, but replaces
the exponential-memory node with a fractional one and the spectral-radius
edge-of-chaos knob with the Matignon wedge. The open theory items are a
fractional ESP theorem and a capacity-redistribution theorem for power-law
kernels.

\section{Neural mass and neural field models}
\label{sec:masses}

\subsection{Classical models}

The biological substrate of the phantom brain comes from two modelling
traditions. \emph{Neural masses} coarse-grain a cortical column into coupled
excitatory/inhibitory populations: Wilson \& Cowan's founding model
\cite{wilson1972excitatory,wilson1973mathematical} with its sigmoidal
response and E/I loop producing multistability and oscillations; the
thalamocortical alpha-rhythm loop of Lopes da Silva et al.\
\cite{lopesdasilva1974model}; and the Jansen--Rit column, the standard
EEG/VEP generator \cite{jansen1995electroencephalogram}. \emph{Neural fields}
replace the lumped populations by a continuum sheet with distance-dependent
connectivity: Amari's analysis of lateral-inhibition (Mexican-hat) fields,
proving existence and stability of localised bump solutions
\cite{amari1977dynamics}, systematised in modern reviews
\cite{coombes2005waves,coombes2010large,ermentrout1998neural} and situated in
the model hierarchy by Deco et al.\ \cite{deco2008dynamic}. Pattern formation
via short-range excitation/long-range inhibition is the same
activator--inhibitor (Turing) mechanism as in reaction--diffusion theory
\cite{murray2002mathematical}.

A unifying observation, due to Deco et al.\ \cite{deco2008dynamic}: every rung
of this reduction ladder assumes \emph{exponential} memory kernels
(second-order synaptic filters, exponential adaptation). Fractional calculus
replaces them with power-law kernels at the cost of a single order parameter
--- the minimal generalisation that can express the scale-free temporal
filtering neural tissue actually exhibits.

\subsection{Fractional evidence at the cellular level}

The strongest empirical motivation is cellular. Lundstrom, Higgs, Spain \&
Fairhall showed that neocortical pyramidal neurons perform
\emph{fractional-order differentiation}: spike-frequency adaptation to slowly
varying input follows a fractional derivative, with memory spanning seconds to
hundreds of seconds \cite{lundstrom2008fractional}. Anastasio had earlier
shown the vestibulo-oculomotor integrator behaves as a fractional-order
integrator \cite{anastasio1994fractional}. If single neurons already compute
fractional operators, fractional-order \emph{population} models are not an
abstraction but a coarse-graining of measured physiology.

\subsection{Fractional-order population models}

The fractional generalisation programme is underway but young: Mondal et al.\
analyse a fractional Wilson--Cowan E/I network, showing the derivative order
acts as an extra bifurcation parameter enriching the dynamical repertoire
\cite{mondal2024emergent}; Gonz\'alez-Ram\'irez formulates a Caputo neural
field with lateral-inhibition kernel and derives approximate travelling waves
whose speed and profile the order modulates \cite{gonzalezramirez2022fractional};
and the fractional recurrent-network literature supplies Mittag-Leffler
stability theory for fractional Hopfield and FitzHugh--Nagumo systems
\cite{wang2014stability,brandibur2022stability}. A genuine gap: no
fractional-order Jansen--Rit/EEG neural-mass model appears in the indexed
literature --- the closest items are the fractional field
\cite{gonzalezramirez2022fractional} and fractional RNN work
\cite{wang2014stability}.

\paragraph{Implications for \texttt{fracres}.}
The library implements the fractional Wilson--Cowan reservoir (stacked E/I
masses, separate $\tau_E^\alpha, \tau_I^\alpha$) and the fractional Amari
field (Mexican-hat ring connectivity, nonlinearity inside the convolution
$w * S(u)$). The missing fractional Jansen--Rit/EEG mass model is a
positioning opportunity for the EEG-fitting phase of the programme.

\section{Critical brain dynamics}
\label{sec:criticality}

\subsection{Neuronal avalanches}

Beggs \& Plenz reported that spontaneous LFP events in cortical slice cultures
distribute in size as a power law $P(S) \sim S^{-\tau}$ with
$\tau \approx 1.5$, and in duration with $\alpha \approx 2.0$ --- the
mean-field exponents of a critical branching process (branching ratio
$\sigma = 1$) \cite{beggs2003neuronal}; the patterns are diverse, precise, and
reproducible over hours \cite{beggs2004neuronal}. The exponents obey the
crackling-noise relation $\gamma = (\alpha - 1)/(\tau - 1) \approx 2$, a
non-trivial consistency check on the critical branching picture
\cite{beggs2003neuronal}. Reviews have elevated avalanches to a candidate
organising principle of cortical activity \cite{plenz2007organizing} and
catalogued the functional benefits of the critical point --- maximal dynamic
range, information capacity, and an optimal storage/transfer trade-off
\cite{shew2013functional,chialvo2010emergent}. Spike-avalanche analyses in
vivo across species and states find near-universal exponents with
state-dependent deviations \cite{ribeiro2010spike,dehghani2012avalanche}.

\subsection{Scepticism and methodology}

The criticality hypothesis has serious critics. Beggs \& Timme catalogue the
failure modes: alternative mechanisms mimicking power laws, thresholding
sensitivity, and the rarity of proper statistical testing
\cite{beggs2012being}. The methodological standard is Clauset, Shalizi \&
Newman: maximum-likelihood exponents with $x_{\min}$ selection, bootstrapped
goodness-of-fit, and likelihood-ratio comparison against log-normal and
exponential alternatives \cite{clauset2009power} --- the estimator
\texttt{fracres} implements for its avalanche exponents. A second pitfall is
\emph{subsampling}: any electrode array records a small fraction of the
network, biasing measured distributions \cite{priesemann2009subsampling}; with
subsampling-invariant analysis, in-vivo spike avalanches look
\emph{slightly subcritical with reverberation} rather than exactly critical
\cite{priesemann2014spike}.

\subsection{SOC, quasi-SOC, and homeostasis}

Self-organised criticality explains scale invariance without parameter
tuning \cite{bak1987self,jensen1998self}, but its standard models require
local conservation --- which neural networks, creating and destroying
activity, violate. Bonachela \& Mu\~noz proved that feedback-driven
non-conservative systems self-organise only to \emph{apparent}, finite-range
scale invariance: quasi-criticality \cite{bonachela2009self}.
Williams-Garc\'ia et al.\ located driven, resource-limited branching models on
a nonequilibrium Widom line of near-critical susceptibility maxima, producing
approximate power laws with near-canonical exponents without exact tuning
\cite{williamsgarcia2014quasicritical}. The biological mechanism that could
hold a network in this regime is homeostatic plasticity --- synaptic scaling
and intrinsic-excitability regulation stabilising average firing rates
\cite{turrigiano2004homeostatic}. The modern consensus picture is therefore
\emph{quasi-criticality maintained by homeostatic control}, not strict SOC.

\subsection{E/I balance}

Excitation/inhibition balance is the empirically identified control
parameter: pharmacological manipulation of inhibition moves avalanche duration
exponents and waiting-time clustering while $\tau$ stays near $1.5$
\cite{lombardi2012balance}; avalanches and long-range temporal correlations
co-emerge in the same human MEG networks and covary with E/I-related state
\cite{poil2012critical,dehghani2012avalanche}. This motivates both the E/I
structure of the Wilson--Cowan reservoir and the qSOC homeostat's placement
as a threshold on population activity.

\paragraph{Implications for \texttt{fracres}.}
The qSOC reservoir implements exactly the quasi-criticality-with-homeostasis
picture: a windowed-energy feedback pulling the network toward a critical set
point \cite{bonachela2009self,turrigiano2004homeostatic}. The library's first
criticality study found power-law-like avalanches ($\tau \approx 1.2$,
duration exponents rising with the fractional order) whose statistics
\emph{survive} a strong perturbation only under homeostatic control ---
consistent with the homeostat's role being robustness of the operating point
rather than the source of the exponents. The methodological lessons of
\cite{clauset2009power,priesemann2014spike} (MLE exponents, subsampling
awareness) are built into \texttt{metrics.py}.

\section{Long-range dependence: theory and estimation}
\label{sec:lrd}

\subsection{Self-similar processes and the Hurst phenomenon}

Empirical LRD began with Hurst's rescaled-range analysis of Nile minima: $R/S
\sim n^{K}$ with $K \approx 0.73 \neq 1/2$ \cite{hurst1951long}. Mandelbrot \&
Wallis separated the phenomenon into the \emph{Joseph effect} (temporal
persistence) and the \emph{Noah effect} (heavy-tailed amplitudes)
\cite{mandelbrot1969noah}, and Mandelbrot \& Van Ness supplied the canonical
model: fractional Brownian motion $B_H(t)$, the self-similar Gaussian process
with stationary increments, and its increment process fractional Gaussian
noise (fGn), whose autocorrelation is non-summable for $H > 1/2$
\cite{mandelbrot1968fractional}. The fGn autocovariance
\begin{equation}
  r(k) = \tfrac{1}{2}\bigl(|k+1|^{2H} - 2|k|^{2H} + |k-1|^{2H}\bigr),
  \qquad r(0) = 1,
  \label{eq:fgn}
\end{equation}
is exactly what the Davies--Harte circulant-embedding algorithm synthesises in
$O(n \log n)$ \cite{davies1987tests} --- the drive generator in
\texttt{fracres} (\texttt{drivers.py}), scaled so the realised process has
unit variance in expectation and the exact $r(k)$. ARFIMA/fractional
differencing $(1-B)^d$ provides the discrete-time sibling with $\beta = 2d$
\cite{granger1980introduction}.

\subsection{Estimators and their failure modes}

Second-order LRD theory defines long memory equivalently by non-summable
covariance, a spectral pole $f(\lambda) \sim |\lambda|^{-2d}$, or aggregated
variance $\mathrm{Var}(X^{(m)}) \sim m^{2H-2}$ \cite{beran1994statistics}.
The estimator landscape is treacherous: Taqqu, Teverovsky \& Willinger's
Monte Carlo study showed most $H$ estimators are strongly biased at realistic
lengths, so agreement among methods is required before claiming LRD
\cite{taqqu1995estimators}. The two estimators in \texttt{fracres} cover the
time and frequency domains: detrended fluctuation analysis
\cite{peng1994mosaic}, for which the DFA exponent equals $H$ for stationary
fGn (and $H+1$ for fBm, whence a stationarity check), and spectral-slope
estimation, since the fGn PSD decays as
\begin{equation}
  S(f) \propto f^{-\beta}, \qquad \beta = 2H - 1
  \label{eq:psd}
\end{equation}
over the scaling range --- so $1/f$-like spectra ($\beta \approx 1$) sit at
the edge of the fGn family. Wavelet log-scale diagrams estimate $H$ from the
scale-energy law $\mathbb{E}|d_{j,k}|^2 \propto 2^{j(2H)}$, exploiting the
decorrelation of wavelet coefficients for fGn-like processes
\cite{abry1998wavelet,veitch1999wavelet}; the wavelet lens connects
self-similarity to Besov regularity \cite{abry2003self}, the function-space
framing the library's regulariser adopts.

\subsection{Heavy tails: the variance trap}

All of the above presupposes finite variance. For $\alpha$-stable processes
with tail index $\alpha_{\mathrm{tail}} < 2$, variance is infinite,
covariance undefined, and \emph{every} $L^2$-based estimator --- R/S,
variance--time, aggregated variance, DFA to second order --- is biased or
meaningless: heavy tails masquerade as, and corrupt estimates of, long memory
\cite{samorodnitsky1994stable,samorodnitsky2007long}. Samorodnitsky's survey
shows LRD is \emph{plural} under heavy tails: scaling of sums, maxima, and
ergodic structure need not agree, and dependence must be characterised by
moment-free measures \cite{samorodnitsky2007long}. The principled response is
to move from covariance to \emph{scale-space regularity}: $p$-th order
wavelet/Besov energies with $p < \alpha_{\mathrm{tail}}$, where the
smoothness index must respect $s < \min\{H, 1/p\}$ --- the $1/p$ ceiling is
the jump/heavy-tail guardrail. This is exactly the Besov index convention
implemented in \texttt{fracres}' \texttt{training.besov\_indices} and its
Littlewood--Paley regulariser, and it frames the library's thesis-level
claim: covariance-LRD is the $p=2$, stationary shadow of a broader
Besov--wavelet scaling phenomenon.

\subsection{LRD in neural signals}

The neural evidence is strong and old: amplitude envelopes of spontaneous
human MEG/EEG oscillations show power-law correlations over hundreds of
seconds with DFA exponents $\approx 0.6$--$0.8$
\cite{linkenkaer2001long}; $1/f$ aperiodic spectra are ubiquitous across
recording modalities and covary with E/I balance, state, and age
\cite{he2014scale,voytek2015dynamic}; and DFA methodology --- including the
$\alpha = H$ vs.\ $H+1$ ambiguity, crossovers, and nonstationarity pitfalls
--- is codified for neurophysiology by Hardstone et al.\
\cite{hardstone2012detrended}, who also survey DFA changes in sleep,
development, and pathology. These are precisely the statistics a phantom
brain must reproduce, and the estimators it must calibrate against.

\paragraph{Implications for \texttt{fracres}.}
The library's metric suite (\texttt{metrics.py}) implements DFA and the
log-binned low-frequency spectral fit with the finite-sample safeguards the
literature demands \cite{taqqu1995estimators,hardstone2012detrended}; the
drive is exact Davies--Harte fGn \cite{davies1987tests}; and the Besov
regulariser enforces the heavy-tail-aware bounds $p < \alpha_{\mathrm{tail}}$,
$s < \min\{H,1/p\}$ \cite{samorodnitsky2007long,abry2003self}. The remaining
gap is a non-Gaussian (stable) drive so the $\alpha_{\mathrm{tail}} < 2$
machinery is actually exercised.

\section{Synthesis: the \texttt{fracres} programme in context}
\label{sec:synthesis}

Table~\ref{tab:map} summarises how each body of theory lands in code. Three
conclusions emerge from the assembled literature.

\begin{table}[t]
\centering\small
\caption{Theory-to-code map for \texttt{fracres}, with the anchoring literature.}
\label{tab:map}
\begin{tabular}{@{}lll@{}}
\toprule
Component & Theory & Key references \\
\midrule
\texttt{kernels.py} & GL / L1-Caputo schemes & \cite{oldham1974fractional,linxu2007l1,scherer2011grunwald} \\
\texttt{validation.py} & $D^\alpha t^\beta$, Mittag-Leffler & \cite{podlubny1999fractional,gorenflo2014mittag} \\
\texttt{stability.py} & Matignon wedge, edge of chaos & \cite{matignon1996stability,bertschinger2004edge} \\
\texttt{reservoirs.py} & ESN split, fractional Wilson--Cowan/Amari & \cite{jaeger2001echo,wilson1972excitatory,amari1977dynamics} \\
qSOC homeostat & Quasi-criticality + homeostasis & \cite{bonachela2009self,turrigiano2004homeostatic} \\
\texttt{drivers.py} & fGn, Davies--Harte synthesis & \cite{mandelbrot1968fractional,davies1987tests} \\
\texttt{regularizers.py} & Besov / Littlewood--Paley & \cite{abry2003self,samorodnitsky2007long} \\
\texttt{training.py} & Readout-only ridge/gradient & \cite{jaeger2001echo,verstraeten2007unification} \\
\texttt{metrics.py} & DFA, spectral slope, CSN MLE & \cite{peng1994mosaic,clauset2009power} \\
\bottomrule
\end{tabular}
\end{table}

\paragraph{1. The niche is genuinely open.}
Fractional RC exists but is four papers deep and none of it targets critical
statistics, neural-mass structure, or stochastic fGn drive
\cite{yao2020fractional,goswami2026longmemory,mastin2026fractional,teuscher2026hardware}.
Fractional neural models exist but no fractional Jansen--Rit/EEG mass model
is indexed (Section~\ref{sec:masses}). The intersection --- fractional-node
reservoirs as \emph{critical brain phantoms} --- appears unoccupied.

\paragraph{2. The biological and theoretical motivations are independent and
convergent.}
Cellular fractional computation \cite{lundstrom2008fractional,anastasio1994fractional},
power-law population filtering (Section~\ref{sec:masses}), quasi-criticality
with homeostasis \cite{bonachela2009self,turrigiano2004homeostatic}, and the
failure of $L^2$ LRD theory under heavy tails
\cite{samorodnitsky2007long} each independently push toward the same
construction: power-law memory operators, criticality control, and
Besov-geometry regularisation.

\paragraph{3. The sharpest open problems are theoretical, and they are named.}
(a)~A fractional echo-state property: contractivity of the fractional state
map under algebraic (non-geometric) decay, in the spirit of
\cite{buehner2006tighter,yildiz2012echo}.
(b)~Memory-capacity redistribution: how power-law kernels reallocate the
Dambre capacity budget across delays \cite{dambre2012capacity}.
(c)~Criticality claims must be made with Clauset--Shalizi--Newman rigour and
subsampling awareness \cite{clauset2009power,priesemann2014spike} ---
avalanches from a phantom are hypotheses about mechanism, not evidence about
brains, until confronted with data.
(d)~The inverse problem: inferring $(\alpha, H)$ and the E/I operating point
from real EEG, where the estimator pathologies of
Section~\ref{sec:lrd} are themselves part of the question.

\appendix

\section{Bibliographic verification}
\label{app:verification}

All 81 references were collected and verified programmatically (September
2026) against live bibliographic APIs --- Crossref/doi.org content
negotiation, OpenAlex, OpenLibrary, NeurIPS proceedings, and arXiv abstract
pages --- with per-entry verification logs in \texttt{docs/litreview/notes/}.
Residual caveats: Matignon 1996 \cite{matignon1996stability} has no DOI
(OpenAlex-verified; page numbers from secondary citations); the Jaeger GMD
reports \cite{jaeger2001echo,jaeger2002memory} are verified via author-hosted
copies; the Abry--Flandrin--Taqqu--Veitch chapter \cite{abry2003self} is
verified at container level only; Mandelbrot \& Wallis is cited per verified
publisher metadata (WRR \textbf{4}(5), 1968), which corrects the widespread
miscitation as vol.~5, 1969 \cite{mandelbrot1969noah}; Voytek \& Knight 2015
appears in \emph{Biological Psychiatry}, not \emph{Neuron}
\cite{voytek2015dynamic}.

\bibliographystyle{plainnat}
\bibliography{references}

\end{document}
